\title{No Cash Flows, New Owners:\\ Spot Bitcoin ETFs and the Origins of Co-movement}

\begin{abstract}
\noindent On January 10, 2024 the SEC approved US spot Bitcoin exchange-traded products, widening who can conveniently hold an asset that has no cash flows. We ask whether this change in the investor base altered Bitcoin's co-movement with US equities. We compare Bitcoin with nine cryptocurrencies that never received such a product, using monthly correlations with the S\&P 500 from 2021 and randomization inference suited to a single treated asset. Bitcoin's equity correlation shows no detectable change relative to the controls: the estimated shift is $-0.0181$ in correlation units, under two correlation points and of the opposite sign to the broadened-access prediction.
% src: estimation_results.json#main.effect_in_correlation_units
The design could detect only an increase of about $0.148$ or larger, the conservative of two power calculations, so an increase of half that size remains possible.
% src: estimation_results.json#power_analysis.mde_correlation_units_ri
Bitcoin's volatility relative to equities fell after the listing. The evidence bounds the part of Bitcoin's 2024 co-movement with equities that the listing can explain.

\medskip
\noindent\textbf{Keywords:} co-movement, exchange-traded funds, investor clienteles, Bitcoin, difference-in-differences

\noindent\textbf{JEL:} G12, G14, G23
\end{abstract}
